\documentclass[11pt,a4paper]{article}

\usepackage[margin=25mm,headheight=14pt]{geometry}
\usepackage{fontspec}
\usepackage{graphicx}
\usepackage{microtype}
\usepackage{booktabs}
\usepackage{siunitx}
\usepackage{tabularx}
\usepackage{xltabular}
\usepackage[inline]{enumitem}
\usepackage{xcolor}
\usepackage{titlesec}
\usepackage{fancyhdr}
\usepackage{caption}
\usepackage{subcaption}
\usepackage{tikz}
\usetikzlibrary{positioning,arrows.meta,fit}
\usepackage[numbers,sort&compress]{natbib}
\usepackage[hidelinks]{hyperref}
\usepackage{cleveref}

\hypersetup{
  pdftitle={POLAR WINDOW: a lifecycle planner for CLPS lunar deliveries},
  pdfsubject={NASA Space Apps Challenge 2026, CLPS Lunar Mission Browser},
  pdfauthor={POLAR WINDOW team},
  pdfkeywords={CLPS, lunar south pole, trajectory, launch weather, radiation, uncertainty},
  pdfcreator={LaTeX}
}

% Palette of the figures, also used for the section rules.
\definecolor{pwblue}{HTML}{2563EB}
\definecolor{pworange}{HTML}{E8710A}
\definecolor{pwgreen}{HTML}{059669}
\definecolor{pwgrey}{HTML}{6B7280}
\definecolor{pwtext}{HTML}{111827}

\titleformat{\section}
  {\normalfont\large\bfseries}{\thesection}{0.6em}{}
  [{\color{pwblue}\titlerule[0.8pt]}]
\titlespacing*{\section}{0pt}{1.4\baselineskip}{0.5\baselineskip}
\titleformat{\subsection}{\normalfont\normalsize\bfseries}{\thesubsection}{0.6em}{}
\titlespacing*{\subsection}{0pt}{0.9\baselineskip}{0.3\baselineskip}

\pagestyle{fancy}
\fancyhf{}
\fancyhead[L]{\footnotesize NASA Space Apps Challenge 2026: CLPS Lunar Mission Browser}
\fancyhead[R]{\footnotesize\thepage}
\fancyfoot[C]{\footnotesize\url{https://github.com/MdAhbab/PolarWindow}}
\renewcommand{\headrulewidth}{0.4pt}
\renewcommand{\footrulewidth}{0.4pt}

\captionsetup{font=small,labelfont=bf,skip=4pt}

\sisetup{
  group-separator={,},
  group-minimum-digits=4,
  per-mode=symbol,
  range-units=single,
  range-phrase={~to~},
  list-units=single,
  uncertainty-mode=separate
}

\setlist[itemize]{noitemsep,topsep=3pt,leftmargin=1.2em}
\setlist[enumerate]{noitemsep,topsep=3pt,leftmargin=1.5em}

\newcolumntype{Y}{>{\raggedright\arraybackslash}X}

\title{\textbf{POLAR WINDOW}\\[3pt]
  \large A lifecycle planner for CLPS lunar deliveries, from launch to return,\\
  with the uncertainty and source of every number}
\author{}
\date{}

\begin{document}
\maketitle
\thispagestyle{fancy}
\vspace{-2.4\baselineskip}
\begin{center}\small
Submission report for the 2026 challenge \emph{CLPS Lunar Mission Browser}.
All numbers were measured on commit \texttt{f421298} of branch \texttt{featured}.
\end{center}

% ---------------------------------------------------------------------------
\section{Introduction}
\label{sec:intro}

NASA's Commercial Lunar Payload Services (CLPS) programme buys the delivery of science
payloads to the Moon from commercial providers \citep{nasaclps}. Four CLPS landers have
flown: Peregrine Mission One, IM-1, Blue Ghost Mission~1 and IM-2. Their reconstructed
trajectories are public in NASA's SPICE archive \citep{naifclps}. The 2026 NASA Space
Apps challenge \emph{CLPS Lunar Mission Browser} asks for a tool that lets planners,
educators and the public compare landing sites and dates, with the Sun and Earth shown
against the local horizon \citep{spaceapps2026}.

The first version of POLAR WINDOW answered one question. At a south-polar site and over
chosen dates, it reported how long the Sun and Earth stay above the site's own skyline.
A delivery raises many more questions. A project manager must know when the vehicle can
leave the pad and whether the weather rules allow it. The trip passes through sunlight,
shadow and the radiation belts. After landing, the team needs the hours in which a
ground station can hear the lander. The launch vehicle's booster comes back to a landing zone on
land or at sea. Each payload must also be built, tested and delivered on
time. Each answer comes from a different tool and data source, and each tool states its
assumptions in its own way, if at all.

This report describes the extension of POLAR WINDOW into a planner that follows one CLPS
delivery through six sections: \emph{Ready}, \emph{Launch}, \emph{Trip}, \emph{Surface},
\emph{Contact} and \emph{Return}. The work makes four contributions.
\begin{enumerate}
\item A single contract for every number. Each value states its provenance (computed,
  published, entered by a named person, or assumed), a margin (a $\pm$, a range or a
  multiplicative factor) or the reason it has none, and the age of any live input. No
  section returns a verdict, a readiness percentage or a combined score.
\item Trajectory, launch, radiation and entry engines checked against flown missions.
  The references are NASA's reconstructions of the four CLPS trips and the measured
  landing of the OSIRIS-REx sample capsule.
\item A process design that separates network access, heavy computation and request
  handling. Every stored answer records the inputs and engine versions it came from.
\item A disclosure of every planned acceptance tolerance that the engines do not meet,
  with the measured value beside it.
\end{enumerate}

The evaluation is bounded by public data. Only four CLPS trips have reconstructions,
and none of the planned missions publishes a trajectory, a radio link budget or a
return state. The engines are therefore checked on the flown cases, and the planned
missions receive designed answers labelled as such. Two external references named in
the plan could not be run: the General Mission Analysis Tool (GMAT) and the SPENVIS web
service. \Cref{sec:eval} states what replaced each of them.

% ---------------------------------------------------------------------------
\section{Related work}
\label{sec:related}

\paragraph{Illumination and Earth visibility at the lunar poles.} Polar lighting has been
mapped from Lunar Orbiter Laser Altimeter (LOLA) topography \citep{smith2010} by casting
horizons and averaging Sun visibility over years \citep{mazarico2011}. Later work located
persistently illuminated regions in camera images \citep{speyerer2013} and in
higher-resolution terrain models \citep{glaser2014}. These studies produce fractions over
long periods on a fixed grid. They answer where the light is on average, not what
happens at one landing point over one mission's dates. They also report no timing
uncertainty per crossing. The surface engine of POLAR WINDOW computes the Sun and Earth
from the Navigation and Ancillary Information Facility's SPICE system
\citep{acton1996,acton2018} with the DE440 ephemeris \citep{park2021}. It reports
each duration with a margin derived from the terrain's vertical error.

\paragraph{Trajectory and launch design.} High-fidelity mission design is available in
open tools such as GMAT, which has a documented verification programme
\citep{hughes2014}. Its building blocks are standard: Lambert solvers for the initial
transfer \citep{izzo2015}, numerical integration under point-mass and harmonic gravity
\citep{montenbruck2000}, and lunar gravity fields from GRAIL \citep{zuber2013}. These
tools require a scripted mission and an analyst. They do not read NASA's CLPS
reconstructions, they do not relate a launch day to the lighting rule at the landing
site, and they do not label an input as published or assumed. POLAR WINDOW does not
compete with them on fidelity. It connects a small, checked subset of their methods to
the landing site, the pad and the reconstructions.

\paragraph{Launch weather and forecast verification.} Launch weather criteria are
published per vehicle as rule sets, for example the Falcon~9 limits on wind at
\qty{162}{ft} and on lightning near the pad \citep{nasa2020wx}. Ensemble systems such as
the Global Ensemble Forecast System (GEFS) give a sample of possible weather
\citep{guan2022}. Lightning is observed from orbit by the Geostationary Lightning Mapper
(GLM) \citep{goodman2013,virts2025}, and long wind records come from the MERRA-2
reanalysis \citep{gelaro2017}. The standard score for a probability forecast is the
Brier score \citep{brier1950}. Public presentations of launch weather usually give one
chance of a rule being broken. They rarely state the sample behind it or whether that
chance has been checked against observations.

\paragraph{Radiation screening.} Trapped belt fluxes are commonly taken from the AE8 and
AP8 models \citep{vette1991,sawyer1976}, implemented in the IRBEM library
\citep{irbem2022} and exposed in Python by SpacePy \citep{niehof2022}. Galactic cosmic
rays (GCR) can be modelled with one solar modulation parameter \citep{matthia2013}, and
measured doses exist in lunar orbit from CRaTER \citep{schwadron2012} and on the far-side
surface from Chang'e-4 \citep{zhang2020}. Large solar particle events are described by
published spectra \citep{king1974}. Engineering tools such as OLTARIS compute dose behind
shielding once a trajectory and a geometry are known \citep{singleterry2011}. They need
an account and a finished design, so they do not serve early screening of several
candidate trips.

\paragraph{Contact planning.} The Deep Space Network (DSN) publishes each antenna's
horizon mask and link parameters \citep{jpl810005}. Delay-tolerant networks route messages
over scheduled contacts with contact graph routing \citep{araniti2015}. Neither source
combines a station's mask with the terrain around a polar lander, which decides whether
Earth is visible at all.

\paragraph{Gap.} Each area has mature tools, but each answers one phase with its own
conventions. None carries the provenance and uncertainty of a number from the pad to the
landing site and back. POLAR WINDOW adds that thread. Its engines are deliberately
simpler than the specialist tools above, and every answer names the specialist tool in
which it should be confirmed.

\bibliographystyle{unsrtnat}
\bibliography{ref}

\end{document}
